\section{Introducción}

\subsection{El problema que no desaparece}
¿Por qué siente usted dolor? Es la pregunta más sencilla del mundo. Y, a la vez, la más difícil. Todos sabemos qué es el dolor. Todos lo hemos sentido. Pero ¿por qué lo sentimos? Sabemos que los nociceptores envían señales al tálamo, que la corteza las procesa y que se liberan endorfinas. Todo eso es una descripción funcional. Nos dice qué hace el cerebro, pero no nos dice por qué usted lo siente.

Piense en un ordenador. Puede procesar información, reconocer una imagen, jugar al ajedrez. Pero no siente nada. No tiene experiencia subjetiva. Nada le duele. ¿Por qué el cerebro es diferente? Este es el problema difícil de la conciencia, formulado por David Chalmers en 1995. Desde entonces, la neurociencia ha hecho enormes progresos, pero el problema difícil ha permanecido intacto. Nadie ha logrado explicar por qué la materia se convierte en experiencia. En 1998 Christof Koch apostó con Chalmers que un correlato neuronal de la conciencia sería identificado en 25 años. En 2023 reconoció que la apuesta estaba perdida.

\subsection{La materia oscura --- las cifras exactas}
El universo tiene una composición que conocemos bastante bien. La misión Planck (2018) nos dio las cifras exactas. Solo el 4,9\% del universo es materia ordinaria --- átomos, estrellas, planetas, personas. La materia oscura representa el 26,8\% y la energía oscura el 68,3\%. En la práctica, el 95,1\% del universo es desconocido.

Imagínese el universo como un cofre con 100 monedas. Solo 5 monedas son visibles. Las otras 95 son invisibles. De ellas, 27 son materia oscura y 68 energía oscura. Sabemos de la materia oscura que existe (las pruebas gravitatorias son sólidas), que no emite ni absorbe luz, y que nos atraviesa sin interactuar. Unos 100 millones de partículas atraviesan su cuerpo cada segundo. No sabemos de qué está hecha.

Una restricción reciente (2025), debida a Bonvin et al., ha demostrado que la materia oscura cae en los pozos gravitatorios exactamente como la materia ordinaria. Si existe una «quinta fuerza», no puede superar el 7\% de la fuerza gravitatoria.

\subsection{¿Qué es ART-12D?}
ART-12D parte de una idea sencilla: el universo tiene más dimensiones que las cuatro que percibimos. Estas dimensiones adicionales son compactas --- tan pequeñas que no las notamos. La materia oscura existe en las doce dimensiones. Proyectada en 4D, aparece como un campo que llena el universo. Este campo es lo que el cerebro recibe como conciencia.

El nombre viene de: \textbf{Akashic} (el campo informacional universal), \textbf{Resonance} (porque la interacción es una resonancia), \textbf{Theory} (un marco teórico) y \textbf{12D} (porque trabajo en doce dimensiones).

Los tres postulados son: 1. La materia oscura es un campo cuántico fundamental. 2. El cerebro es un receptor geométrico. 3. El acoplamiento se produce por interacción de Yukawa.

Piense en una radio. No crea la música, la capta del aire. Las ondas de radio existen tanto si la radio está encendida como si no. Cuando la radio se estropea, la música no desaparece del universo --- solo cesa su recepción local. Del mismo modo, el cerebro funciona exactamente como una radio. La conciencia no la generan las neuronas, la reciben. Y lo que reciben es el campo de la materia oscura.

\subsection{¿Por qué doce dimensiones?}
¿Por qué doce? ¿Por qué no 10, o 11, o 26? La respuesta viene de la teoría de cuerdas. Imagine una cuerda de guitarra. Puede vibrar de varias maneras. Cada manera corresponde a una nota. Del mismo modo, una cuerda fundamental puede vibrar en varias dimensiones. El número de dimensiones determina qué tipos de vibración son posibles.

La historia de las dimensiones: Kaluza--Klein (5 dimensiones, 1919--1926), supercuerdas (10, 1970--1980), teoría M (11, Witten, 1995), teoría F (12, Vafa, 1996), ART-12D (12, Ofiteru, 2026). Doce es el mayor número de dimensiones en el que la teoría de cuerdas es matemáticamente consistente. Más dimensiones llevarían a inconsistencias; menos, a una teoría incompleta.

\subsection{¿Cómo son las dimensiones compactas?}
Imagine una manguera de jardín. De lejos parece una línea 1D --- solo tiene longitud. De cerca tiene grosor --- es 2D (la superficie del cilindro). Del mismo modo, nuestro espacio parece tridimensional. Pero de cerca (en la escala de Planck, $10^{-35}$ m) podría tener once dimensiones. Las ocho dimensiones compactas están enrolladas en cada punto del espacio.

¿Qué significa «enrolladas»? No son lugares a los que se viaja. No son mundos paralelos a los que se llega viajando. Son propiedades locales de cada punto. Cada punto del espacio-tiempo tiene, además, ocho pequeñas dimensiones «unidas» a él.

\subsection{¿Cómo se relacionan las dimensiones con la conciencia?}
La materia oscura existe en doce dimensiones. Es un campo escalar $A(x,y)$, donde $x$ son las coordenadas 4D e $y$ las ocho coordenadas compactas. Al compactificar las ocho dimensiones, el campo se descompone en modos de Kaluza--Klein:
\begin{equation}
A(x,y) = \sum_n A_n(x) Y_n(y)
\end{equation}
Cada modo tiene una masa y una frecuencia característica:
\begin{equation}
m_n^2 = \frac{n^2}{R^2}
\end{equation}
El cerebro, por su estructura geométrica (los microtúbulos de las neuronas), puede resonar con ciertos modos. Cuando resuena, el campo se transduce en experiencia subjetiva.

¿Por qué importa la arquitectura y no el tamaño? El cerebelo tiene 69.000 millones de neuronas (el 80\% del total), pero no produce conciencia. Tiene una arquitectura \textit{feedforward} --- la información fluye en una sola dirección, como una antena de hilo. La corteza tiene solo 16.000 millones de neuronas, pero es la sede de la conciencia. Tiene una arquitectura \textit{recurrente} --- la información circula en bucles, como una antena parabólica bien orientada. Conclusión: no cuenta el número de neuronas, sino la geometría de la red.

\section{Metodología}
\subsection{Cómo trabajé}
No soy un físico profesional. No tengo laboratorio. No tengo financiación. Lo que tengo es una biblioteca, tiempo y una pregunta que no me deja en paz. Trabajé en cuatro pasos. El primer paso fue la formulación de la hipótesis, partiendo de la intuición de que la conciencia se recibe y no se genera. El segundo paso fue la formalización, con las herramientas de las matemáticas conocidas: Kaluza--Klein para las dimensiones adicionales, la teoría F para doce dimensiones, la teoría de campos para los lagrangianos, el IIT para la información integrada. El tercer paso fue la deducción de predicciones, es decir, lo que deberíamos observar si la teoría es correcta. Seleccioné las predicciones testables hoy. El cuarto paso fue la propuesta de experimentos con medios modestos, empezando por uno de 50.000 dólares.

\subsection{Qué me guio}
Me até a cuatro criterios. La coherencia interna --- las ecuaciones no deben contradecirse. La falsabilidad --- la teoría debe hacer predicciones testables. La convergencia --- la teoría no debe contradecir la física establecida. La honestidad --- debo reconocer que es especulativa.

\subsection{Qué utilicé}
Utilicé Kaluza--Klein (1919) para la idea de dimensiones compactas, la teoría F (Vafa, 1996) para el marco de doce dimensiones, el IIT (Tononi, 2004) para la medida de la información integrada, Orch OR (Penrose, 1996) para el colapso cuántico en los microtúbulos, y la teoría cuántica de campos para los lagrangianos y las ecuaciones de movimiento.

\subsection{Mis limitaciones}
Debo ser honesto sobre las limitaciones. No tengo pruebas experimentales; la teoría no ha sido testada. No tengo revisión por pares; no he publicado en revistas científicas. No tengo afiliación académica; soy investigador independiente. No tengo financiación; mis medios son limitados.
